\documentclass[12pt]{article}
\usepackage[margin=2.5cm]{geometry}
\usepackage{amsmath,amssymb,booktabs,longtable,graphicx,array,setspace,hyperref,multicol}
\usepackage{marginnote}
\usepackage{xcolor}
\usepackage{caption}
\usepackage{subcaption}
\usepackage{enumitem}
\captionsetup{font=small,labelfont=bf}
\setstretch{1.25}
\newcommand{\Stk}{\mathrm{Stk}}
\newcommand{\Rey}{\mathrm{Re}}
\newcommand{\PiVC}{\Pi_{VC}}

% --- Nomenclature entry format ---
\newcommand{\nomitem}[2]{#1 & #2 \\}

\title{Predicting cyclone collection efficiency across geometries: \\
validation, physics-inspired feature ablation and uncertainty assessment}
\author{A.~Hanifi$^{a}$, A.~Beniaiche$^{a,*}$, A.~Belkalouche$^{b}$, S.~Ghodbane$^{b}$, C.~Belkacemi$^{b}$, M.~Belkadi$^{c}$\\[2pt]
{\small $^a$Laboratory of Fluid Mechanics, Ecole Militaire Polytechnique, Algiers, Algeria}\\
{\small $^b$Ecole Sup\'erieure des Techniques de l'A\'eronautique, Dar El-Beida, Algiers, Algeria}\\
{\small $^c$Laboratory of Turbomachinery, Ecole Militaire Polytechnique, Algiers, Algeria}\\
{\small $^*$Corresponding author: beniaiche\_ahmed25@yahoo.fr}}
\date{Revised version (R1) of POWTEC-D-26-03167}

\begin{document}
\maketitle

% =====================================================================
%  ABSTRACT
% =====================================================================
\begin{abstract}
Machine-learning models of cyclone collection efficiency are nearly always validated on random point-wise splits, which distribute data points belonging to the same physical cyclone across both training and test sets. This practice inflates apparent accuracy because the model is evaluated on interpolation along a grade-efficiency curve rather than on genuine extrapolation to an unseen design. We revisit this question using the publicly available 977-point, 28-geometry efficiency database of Bregolin et al.\ (Waste 2, 2024). Three questions guide the work: (i) how much accuracy is actually lost when whole geometries are withheld from training; (ii) whether physics-inspired input features---a vortical-capture group $\PiVC$, a separation number, and an inlet-geometry ratio---improve prediction in a controlled ablation; and (iii) whether the extrapolation and uncertainty diagnostics that one might recommend to practitioners are themselves validated. Gradient-boosting models attain $R^2\approx0.95$ under random splitting but only $R^2\approx0.73$--0.75 (RMSE $\approx 0.17$) when whole geometries are withheld, whether by 20 repeated 23/5 geometry-grouped splits or by leave-one-geometry-out (LOGO) validation; semi-empirical Lapple and Stokes-only baselines give $R^2\approx0.29$ and 0.44 respectively. Six tree-based learners and one neural network are compared under identical protocols. The 17 engineered features yield small, consistent gains under tuned hyperparameters (RMSE reduction 0.001--0.011, at least 7 of 10 splits improved) but no statistically significant gain under LOGO: the paired bootstrap 95\% interval of the macro-RMS difference $\Delta\text{RMSE}_{\text{macro,RMS}}$ contains zero for every model. A large gain is observed only for the neural network under grouped splits with fixed hyperparameters ($R^2$ 0.51$\to$0.69), but this gain does \emph{not} transfer to LOGO ($R^2_{\text{LOGO}}$: 0.598$\to$0.594). The constants of $\PiVC$, calibrated strictly inside training folds, are poorly identified; a freely fitted Reynolds exponent converges to $m\approx+0.43$, not the assumed $-0.5$, reflecting the intrinsic dependence of Stk on Re. A Mahalanobis distance in geometry space does not significantly predict LOGO error (Spearman $\rho=0.13$, $p=0.50$). A heteroscedastic interval model under-covers at the 95\% level (90\% at $\pm2\sigma$ versus 95.4\% nominal); a geometry-calibrated conformal rescaling achieves 92\% pooled and 95\% mean per-geometry coverage but only 64.5\% for the worst geometry---so this represents empirical mean coverage within the available database, not a formal guarantee for an arbitrary unseen geometry. For the cut diameter, the Reynolds exponent is $-0.46$ (pooled), $-0.55$ (geometry-adjusted) and $-0.52$ (fixed effects, wide interval), but Re and geometry are 98.7\% collinear in this database, so the exponent cannot be disentangled from geometry effects. Taken together, these results offer a conservative, reproducible performance baseline and an honest account of what this database can and cannot support.
\end{abstract}

\noindent\textbf{Keywords:} cyclone separator; collection efficiency; gradient boosting; physics-inspired features; geometry-grouped validation; uncertainty quantification

% =====================================================================
%  NOMENCLATURE (after abstract/keywords)
% =====================================================================
\begin{table}[!ht]\centering
\caption{Nomenclature and abbreviations.}\label{tab:nomenclature}
\small
\begin{tabular}{llll}
\toprule
\multicolumn{2}{l}{\textbf{Latin symbols}} & \multicolumn{2}{l}{\textbf{Greek symbols}} \\
\cmidrule(lr){1-2}\cmidrule(lr){3-4}
Symbol & Definition & Symbol & Definition \\
\midrule
$D$ & Cyclone body diameter [m] & $\eta$ & Collection efficiency [--] \\
$D_e$ & Vortex finder diameter [m] & $\alpha$ & Reynolds exponent of Stk$_{50}$ [--] \\
$D_p$ & Particle diameter [m] & $\rho_p$ & Particle density [kg/m$^3$] \\
$a$ & Inlet height [m] & $\rho_g$ & Gas density [kg/m$^3$] \\
$b$ & Inlet width [m] & $\nu$ & Kinematic viscosity [m$^2$/s] \\
$N$ & Number of effective turns [--] & $\sigma$ & Standard deviation of residuals [--] \\
$\Rey$ & Reynolds number, $v_\text{in}D/\nu$ [--] & $\tau_p$ & Particle relaxation time [s] \\
$\Stk$ & Stokes number, $\tau_p v_\text{in}/D$ [--] & \multicolumn{2}{l}{\textbf{Subscripts}} \\
$\Stk_{50}$ & Stokes number at $\eta=0.5$ [--] & $g$ & Per-geometry index \\
$v_\theta$ & Tangential velocity [m/s] & $\text{ref}$ & Reference value \\
$v_r$ & Radial velocity [m/s] & $50$ & At 50\% efficiency \\
$\PiVC$ & Vortical-capture group [--] & & \\
$A_\Pi$ & Calibration constant of $\PiVC$ [--] & & \\
$B_\Pi$ & Damping constant of $\PiVC$ [--] & & \\
$C_\Pi$ & Density-ratio exponent of $\PiVC$ [--] & & \\
$m_\Pi$ & Reynolds exponent of $\PiVC$ [--] & & \\
$G$ & Number of geometries [--] & & \\
\midrule
\multicolumn{4}{l}{\textbf{Abbreviations}} \\
\midrule
LOGO & Leave-One-Geometry-Out & BPNN & Back-Propagation Neural Network \\
L2GO & Leave-Two-Geometries-Out & MLP & Multi-Layer Perceptron \\
GBR & Gradient Boosting Regressor & RF & Random Forest \\
XGB & Extreme Gradient Boosting & HGB & Histogram Gradient Boosting \\
LGBM & Light Gradient Boosting Machine & CV & Coefficient of Variation \\
RMSE & Root Mean Square Error & MAE & Mean Absolute Error \\
AUC & Area Under the (ROC) Curve & CI & Confidence Interval \\
PCHIP & Piecewise Cubic Hermite Interpolation & OOF & Out-of-Fold \\
\bottomrule
\end{tabular}
\end{table}

% =====================================================================
%  INTRODUCTION
% =====================================================================
\section{Introduction}

Cyclone separators are among the most widely deployed gas--solid separation devices in industrial practice, owing to their simplicity, lack of moving parts, and tolerance to high-temperature and high-pressure environments \cite{hoffmann}. The collection efficiency of a given cyclone depends on a confluence of factors---particle inertia, gas flow regime, and the detailed geometry of the inlet, body, and cone---that semi-empirical models such as those of Lapple \cite{lapple}, Barth \cite{barth}, and Muschelknautz \cite{muschel} capture only to a limited degree \cite{hoffmann92,dirgo}. Machine-learning models have been investigated for cyclone collection-efficiency prediction \cite{bregolin}, often reporting coefficients of determination exceeding 0.95. Related studies have also examined the influence of inlet geometry on cyclone flow and performance using CFD \cite{elsayed}, as well as artificial-neural-network prediction of cyclone pressure-drop coefficients \cite{zhao}. These numbers, however, are almost invariably obtained through random point-wise splitting of the dataset: a random subset of measurements is held out for testing, while the remainder is used for training. Because multiple measurements from the same physical cyclone---corresponding to different particle sizes or operating conditions---are distributed across both sets, the model is effectively asked to interpolate along an already-seen grade-efficiency curve rather than to predict the performance of a genuinely new design. The reported accuracy therefore measures something closer to curve fitting than to design generalization.

The distinction matters. A practitioner who consults a published model to estimate the efficiency of a proposed new cyclone geometry needs to know how the model will perform on \emph{unseen} designs, not on additional operating points of designs that were already in the training data. Yet very few studies in the cyclone or broader particulate-flow literature report geometry-level generalization or calibrated uncertainty estimates. The present paper is deliberately framed to fill this gap: it is a validation and ablation study, not a proposal for a new learning algorithm. The learner is a standard tree ensemble (or, for comparison, a back-propagation neural network of the type used by the original database authors); physics enters only through the construction of input features. We do not claim this to be ``physics-informed machine learning'' in the formal sense of Karniadakis et al.\ \cite{karniadakis}, where governing equations constrain the model architecture or loss function.

Our specific contributions are as follows:
\begin{enumerate}[label=(\roman*)]
\item a systematic quantification of the accuracy lost between random and geometry-level validation, using repeated grouped splits, LOGO, and leave-two-geometries-out (L2GO) protocols;
\item a controlled ablation in which identical models, trained under an identical hyperparameter budget, are compared with and without engineered features---with the constants of the physics-inspired group calibrated strictly inside training folds to prevent data leakage;
\item an evaluation of a Mahalanobis extrapolation indicator and of both heteroscedastic and conformal prediction intervals under the LOGO protocol;
\item a cut-diameter analysis with propagated uncertainty, including an explicit treatment of the confounding between Reynolds number and geometry that is intrinsic to this database.
\end{enumerate}

Recent work on hydrocyclones has explored adaptive optimization and control under dynamic particle-size distributions \cite{rineng}, CFD-trained CNN--LSTM surrogates for instantaneous flow-field prediction \cite{mineng,e2024}, hybrid machine-learning/mechanistic models for internal-state prediction \cite{seppur}, and multi-objective inlet optimization \cite{e2025}. These studies, however, typically target flow fields or process control for a single design family and are validated on data from that same family. Our concern is different: we ask how well a performance model transfers \emph{between} different cyclone geometries, and we provide an honest accounting of where current data-driven models succeed and where they fall short.

% =====================================================================
%  DATA
% =====================================================================
\section{Data}

\subsection{Efficiency database}
The dataset used throughout this study is the one published as supplementary material by Bregolin et al.\ \cite{bregolin}. It contains 980 records, of which three are exact duplicates, leaving $N=977$ unique measurements after de-duplication. Each record consists of seven dimensionless inputs---$D_e/D$, $a/D$, $b/D$, $\Rey$, $\Stk$, $D_p/D$, $\rho_p/\rho_g$---and the collection efficiency $\eta$. A \emph{geometry} is defined by the triplet $(D_e/D, a/D, b/D)$; the database contains 28 such geometries, with between 7 and 161 points per geometry. Other design variables that are known to affect cyclone performance---body and cone heights, inlet type, dipleg geometry---are not recorded in the database, so two physically distinct cyclones that happen to share the three ratios are merged into a single geometry identity. A \emph{condition} is defined as one geometry at one Reynolds number and density ratio; the database contains 73 conditions, each contributing one grade-efficiency curve. Table~\ref{tab:ranges} gives the parameter ranges; these replace the incorrect ranges that appeared in the first submission.

\begin{table}[!ht]\centering\caption{Parameter ranges (977 unique points, 28 geometries).}\label{tab:ranges}\input{t_ranges.tex}\end{table}

An important structural feature of the database is the near-collinearity between the Stokes number and the other inputs. The Stokes number is, to a good approximation, a function of the remaining variables:
\begin{equation}
\Stk \approx \frac{\rho_p}{\rho_g}\left(\frac{D_p}{D}\right)^2 \frac{\Rey}{18},
\end{equation}
where the ratio between the actual Stk and this expression lies in the range 1.0--1.1 for 75\% of the points (the small deviations reflecting rounding of $D_p/D$). The input set is therefore strongly collinear, which has consequences for feature engineering: adding features that are deterministic functions of existing inputs cannot, by itself, introduce new information into the model. It can, however, change the inductive bias of the learner by making certain nonlinear relationships more accessible to tree-based splits or neural activations.

The database aggregates experiments from several literature sources, each with its own test rig, particle material, and measurement methodology. Source labels are not included in the released data, so inter-source bias cannot be modelled explicitly. We characterize the heterogeneity of each geometry in Appendix~A (Reynolds range, density ratio, sample size). A particularly striking property is that the Reynolds number is almost a property of the geometry rather than an independent variable: the correlation between $\ln\Rey$ and its geometry-level mean is 0.987 at the condition level. This means that any regression coefficient attributed to Re in a pooled fit partly reflects between-geometry differences, a confounding that we revisit in the cut-diameter analysis (Section~4.7). The pressure-drop data of Schmid and Vogel that were analysed in the first submission are not part of this database and are not considered further here (see Section~\ref{sec:limits}).

% =====================================================================
%  METHODS
% =====================================================================
\section{Methods}

\subsection{Physics-inspired features}
\label{sec:pivc}

The starting point for feature engineering is the classical impaction argument. A particle suspended in a swirling flow experiences, in the Stokes regime, a radial drift driven by the centrifugal acceleration and opposed by viscous drag. Writing the tangential velocity as $v_\theta$ and the particle relaxation time as $\tau_p$, the radial equation of motion is
\begin{equation}
\tau_p\,\frac{\mathrm{d}v_r}{\mathrm{d}t} = \frac{\tau_p\, v_\theta^2}{r} - v_r,
\label{eq:eom}
\end{equation}
which corrects Eq.~(1) of the first submission (the factor $\tau_p$ was previously omitted). Scaling this expression at a characteristic separation radius suggests that capture depends on the Stokes number and on a drag-related correction involving the Reynolds number. We therefore introduce the semi-empirical vortical-capture group
\begin{equation}
\PiVC = 1 - \exp\!\left[-\frac{A_\Pi\,\Stk\,\Rey^{m_\Pi}}{1 + B_\Pi\,(\rho_p/\rho_g / 1000)^{C_\Pi}}\right],
\label{eq:pivc}
\end{equation}
with a nominal exponent $m_\Pi = -1/2$. This $\Rey^{-1/2}$ scaling is an exploratory scaling hypothesis, not derived from the equations presented in this study. The constants $(A_\Pi, B_\Pi, C_\Pi)$ are fitted to data, and a freely fitted Reynolds exponent converges to $m_\Pi \approx +0.43$ rather than $-0.5$, as discussed in Section~4.3. The density-ratio term in the denominator introduces an empirical correction to the drag-related term through the particle-to-gas density contrast, which would otherwise be absent from the Stokes number alone.

The calibration of $(A_\Pi, B_\Pi, C_\Pi)$ is performed by nonlinear least-squares minimization of the objective function
\begin{equation}
J(A_\Pi, B_\Pi, C_\Pi) = \sum_{i \in \mathcal{T}} \left[\eta_i - \PiVC{}_i(A_\Pi, B_\Pi, C_\Pi)\right]^2,
\label{eq:objective}
\end{equation}
where $\mathcal{T}$ denotes the training set of the current fold. This objective operates on the raw efficiency $\eta$ (no weighting, no transformation); it is an unweighted nonlinear least-squares fit of $\PiVC$ to $\eta$. The optimization uses four starting values for $A_\Pi$ ($10^2, 10^3, 10^4, 10^5$), with the solution of lowest cost retained. The bounds are:
\begin{equation}
A_\Pi \in [10^{-3},\, 10^9], \quad B_\Pi \in [0,\, 50], \quad C_\Pi \in [0.1,\, 6], \quad m_\Pi \in [-1.5,\, 0.5] \text{ (when free)}.
\end{equation}
Solutions that reach a bound are retained and reported explicitly (Section~4.3).

To prevent data leakage, $(A_\Pi, B_\Pi, C_\Pi)$ are re-fitted on the training data of each fold only---never on the test fold. Whenever $\PiVC$ appears in the feature set during nested hyperparameter tuning, the constants are re-calibrated inside every inner training fold and applied to the inner validation fold. After hyperparameters are selected, the constants are recalibrated once on the full outer training set, and the outer test geometries are evaluated only once. This nested protocol ensures that the test data never influence the calibration.

Beyond $\PiVC$, the remaining engineered features are: $\Stk/\sqrt{\Rey}$ (a separation number), its logarithm $\ln(\Stk/\sqrt{\Rey})$, the inlet-geometry ratio $ab/D_e^2$, the interaction $\Stk \cdot ab/D_e^2$, the product $\Rey\,D_e/D$, and logarithmic transforms of Re, Stk, $D_p/D$, and $\rho_p/\rho_g$. In total this gives 7 raw inputs + 4 logarithmic transforms + 6 derived features = 17 inputs (the first submission incorrectly stated 14). Values that are non-positive before logarithmic transformation are handled by adding a small epsilon ($10^{-12}$) to avoid numerical issues; no clipping or normalization is applied to the raw inputs, as tree-based models are invariant to monotone transformations. The complete list, with computational names, mathematical definitions, and numerical treatment, is given in Table~\ref{tab:featurelist}.

\begin{table}[!ht]\centering
\caption{Complete feature list (17 inputs). All predictions are clipped to $[0,1]$.}
\label{tab:featurelist}\small
\begin{tabular}{llll}
\toprule
No. & Name & Definition & Treatment \\
\midrule
1 & De & $D_e/D$ & raw \\
2 & a & $a/D$ & raw \\
3 & b & $b/D$ & raw \\
4 & Re & $\Rey$ & raw \\
5 & Stk & $\Stk$ & raw \\
6 & Dp & $D_p/D$ & raw \\
7 & rr & $\rho_p/\rho_g$ & raw \\
\midrule
8 & lnRe & $\ln(\Rey)$ & log, $\epsilon=10^{-12}$ if $\leq0$ \\
9 & lnStk & $\ln(\Stk)$ & log, $\epsilon=10^{-12}$ if $\leq0$ \\
10 & lnDp & $\ln(D_p/D)$ & log, $\epsilon=10^{-12}$ if $\leq0$ \\
11 & lnrr & $\ln(\rho_p/\rho_g)$ & log, $\epsilon=10^{-12}$ if $\leq0$ \\
\midrule
12 & PiVC & Eq.~\eqref{eq:pivc} & calibrated per fold, $m=-0.5$ \\
13 & sep & $\Stk/\sqrt{\Rey}$ & derived \\
14 & lnsep & $\ln(\Stk/\sqrt{\Rey})$ & log of derived \\
15 & geo & $ab/D_e^2$ & derived \\
16 & Stk\_geo & $\Stk \cdot ab/D_e^2$ & interaction \\
17 & Re\_De & $\Rey \cdot D_e/D$ & interaction \\
\bottomrule
\end{tabular}
\end{table}

The five ablation configurations are:
\begin{itemize}[noitemsep]
\item \textbf{Raw} (7): the seven raw inputs only.
\item \textbf{Raw + logs} (11): raw inputs plus the four logarithmic transforms.
\item \textbf{Raw + $\PiVC$} (8): raw inputs plus the vortical-capture group.
\item \textbf{Raw + separation/geometry} (12): raw inputs plus the five derived terms other than $\PiVC$.
\item \textbf{All features} (17): the complete set.
\end{itemize}

\subsection{Models and training protocol}
Six tree-based learners and one neural network are compared under a single, unified protocol. The tree models are: gradient boosting (scikit-learn GBR), XGBoost, LightGBM, CatBoost, random forest, and histogram gradient boosting. The neural network is a BPNN-type architecture (three hidden layers of 15 ReLU units each, Adam optimizer, learning rate 0.008, batch size 5, 100 epochs), matching the pseudo-code distributed with the original database \cite{bregolin}.

Two hyperparameter protocols are used for all boosters and forests:
\begin{itemize}[noitemsep]
\item \textbf{Fixed}: the same hyperparameters are used for the raw and 17-feature inputs of a given model, but the exact values can differ between model families. For the gradient-boosting variants (GBR, XGB, LGBM) and the histogram booster (HGB) we use 400 trees, maximum depth 4, learning rate 0.05 and subsample 0.85; CatBoost uses the same settings without row subsampling; the random forest uses 300 trees with \texttt{min\_samples\_leaf}=2 and no subsampling; and the MLP uses three hidden layers of 15 ReLU units each, Adam with learning rate 0.008, batch size 5 and 100 epochs. This protocol ensures that any performance difference is attributable to the feature set, not to hyperparameter tuning.
\item \textbf{Tuned}: for each outer split, a four-configuration search (depth 3--5, 200--500 trees, learning rate 0.03--0.1; leaf size and feature fraction for RF), scored by 3-fold geometry-grouped inner cross-validation on the training geometries only. The nested protocol operates as follows:
\end{itemize}

\begin{equation}
\begin{aligned}
&\text{Outer training geometries} \\
&\quad\downarrow \\
&\text{Inner training folds (3-fold GroupKFold)} \\
&\quad\downarrow \\
&\text{Calibration of } (A_\Pi, B_\Pi, C_\Pi) \text{ inside each inner training fold} \\
&\quad\downarrow \\
&\text{Transformation applied to inner validation fold} \\
&\quad\downarrow \\
&\text{Selection of hyperparameters (lowest inner MSE)} \\
&\quad\downarrow \\
&\text{Recalibration of } (A_\Pi, B_\Pi, C_\Pi) \text{ on full outer training set} \\
&\quad\downarrow \\
&\text{Evaluation on outer test geometries (single pass)}
\end{aligned}
\label{eq:nested}
\end{equation}

The neural network is not tuned; it uses the architecture and hyperparameters specified above for all experiments. All predictions are clipped to $[0, 1]$.

\subsection{Validation schemes}
Five validation schemes are employed, each answering a different question:
\begin{enumerate}[label=(\alph*)]
\item \textbf{Random point-wise 80/20} (20 repeats): the protocol used in earlier work, in which individual data points are randomly assigned to train or test without regard to geometry. This measures interpolation along known curves.
\item \textbf{Repeated geometry-grouped 23/5} (20 repeats): 23 geometries are randomly selected for training and 5 for testing. This measures transfer to unseen designs.
\item \textbf{LOGO} (28 folds): each geometry is held out in turn. This is the most conservative within-database generalization test.
\item \textbf{L2GO} (60 random pairs): two geometries are held out simultaneously, providing a harder extrapolation test.
\item \textbf{Leave-one-Re-band-out}: an entire Reynolds-number band is withheld, testing extrapolation in operating regime.
\end{enumerate}
Because $R^2$ computed within a geometry that has only 7--10 points is numerically unstable (it can take arbitrarily large negative values when the within-geometry variance is small), we report \emph{pooled out-of-fold} $R^2$, RMSE, and MAE as the primary metrics. The sample-size-weighted mean of per-geometry $R^2_g$ is reported only for comparison with earlier work. Additionally, to address the concern that pooled RMSE gives undue weight to geometries with many points, we report the macro-RMS of per-geometry RMSE:
\begin{equation}
\text{RMSE}_{\text{macro,RMS}} = \sqrt{\frac{1}{G}\sum_{g=1}^{G} \text{RMSE}_g^2},
\end{equation}
which gives equal weight to each geometry regardless of sample size.

\subsection{Semi-empirical baselines}
Two baselines are included to contextualize the machine-learning results. The first is Lapple's grade-efficiency curve,
\begin{equation}
\eta = \frac{1}{1 + (\Stk_{50}/\Stk)^2}, \quad \text{with } \Stk_{50} = \frac{b/D}{4\pi N},
\end{equation}
where $N$ is the number of effective turns (taken as $N=5$ or fitted on the training geometries). The second is a Stokes-only fit, $\eta = 1 - \exp(-A\,\Stk^{k})$, where $A$ and $k$ are fitted on training data. These two semi-empirical models provide reference performance levels using substantially fewer inputs than the machine-learning models. The Stokes-only fit relies on Stk alone, whereas Lapple's model additionally incorporates the inlet-width ratio $b/D$ and the effective number of turns $N$.

\subsection{Extrapolation indicator and prediction intervals}
For each geometry we compute the Mahalanobis distance of $(D_e/D, a/D, b/D)$ from the centroid and covariance of the 27 other geometries (each geometry counted once). The flagging threshold is the 95th percentile of the distances of those 27 geometries. We test the association between this distance and LOGO error by Spearman correlation (with permutation $p$-value) and by the area under the ROC curve (AUC) for identifying the worst-quartile geometries.

Interval half-widths are based on a heteroscedastic model $\sigma_\eta(\hat\eta) = c_2\hat\eta^2 + c_1\hat\eta + c_0$, fitted to binned out-of-fold residuals of the 27 \emph{other} geometries (leave-one-geometry-out). Coverage is then evaluated on the held-out geometry. As an alternative, we use a scale-normalized conformal quantile $q$ of $|\text{res}|/\sigma_\eta$ computed on the other geometries and applied as $\hat\eta \pm q\,\sigma_\eta(\hat\eta)$.

\subsection{Cut diameter}
For each condition, $\Stk_{50}$ is the Stokes number at $\eta = 0.5$, extracted by monotone PCHIP interpolation in $\ln\Stk$ when the data bracket 0.5 (70 of 73 conditions), or by inverting a fitted $\eta = 1 - \exp(-A\,\Stk^{k})$ (3 conditions). Uncertainty is obtained by 1000 bootstrap replicates with Gaussian noise $\sigma_\epsilon = 0.025$ on $\eta$ (assumed, not measured). We fit $\ln\Stk_{50} = \alpha\ln\Rey + \dots$ by weighted least squares in three specifications: (i) Re only; (ii) Re plus geometry and density covariates; (iii) geometry fixed effects. Confidence intervals come from 2000 bootstrap resamples of \emph{geometries} (not conditions), because conditions belonging to the same geometry are not independent.

% =====================================================================
%  RESULTS
% =====================================================================
\section{Results and discussion}

\subsection{Random versus geometry-level validation}
Table~\ref{tab:gap} presents the central finding of this study. Under random point-wise splitting, GBR achieves $R^2 = 0.948 \pm 0.010$, consistent with values reported in earlier work. When whole geometries are withheld, the same model drops to $0.727 \pm 0.143$ over 20 grouped splits and $0.731$ pooled under LOGO. The standard deviation of 0.10--0.14 across grouped splits shows that a single 23/5 split---such as the one used in the first submission---is not a reliable estimate of generalization performance. Table~\ref{tab:fixed} compares all learners under the fixed-hyperparameter protocol.

\begin{table}[!ht]\centering\caption{Accuracy under three validation schemes.}\label{tab:gap}\input{t_gap.tex}\end{table}
\begin{table}[!ht]\centering\caption{20 repeated geometry-grouped 23/5 splits, fixed hyperparameters (mean$\pm$sd of $R^2$; mean RMSE).}\label{tab:fixed}\input{t_fixed.tex}\end{table}

All tree ensembles exhibit similar average performance under grouped validation (differences $<0.03$ against a split-to-split standard deviation exceeding 0.10), so no clear ranking can be established on this database. CatBoost, which appeared highest on a single split in the first submission, is not highest here. All machine-learning models clearly outperform the semi-empirical baselines, which supports the use of multivariable models even though the absolute accuracy for unseen geometries remains moderate.

\subsection{Ablation of the physics-inspired features}
\label{sec:ablation}

Under fixed hyperparameters (Table~\ref{tab:fixed}), the 17-feature set changes $R^2$ for tree ensembles by between $-0.007$ and $+0.015$---within noise---while it improves the neural network from $R^2=0.51$ to $R^2=0.69$. Under tuned hyperparameters (Table~\ref{tab:tuned}), all six tree learners show a small improvement in mean RMSE (0.001--0.011), with at least 7 of 10 splits improved for each learner (10/10 for HGB).

\begin{table}[!ht]\centering\caption{Tuned protocol (10 grouped splits): raw versus 17 features. $\Delta$RMSE is (17 features)$-$(raw); negative is better.}\label{tab:tuned}\input{t_tuned.tex}\end{table}

Under LOGO (Table~\ref{tab:abl}), the picture is more nuanced. Table~\ref{tab:abl} separates the contributions of each feature group for all six tree models. Several observations stand out:
\begin{enumerate}[noitemsep]
\item Logarithmic transforms alone change nothing---the tree models are invariant to monotone transformations of their inputs, so the logarithms add no new splitting information.
\item $\PiVC$ alone changes little (best case HGB: RMSE 0.176$\to$0.172; CatBoost: 0.187$\to$0.184).
\item The largest gains, up to 0.02--0.03 in $R^2$ for XGB, LightGBM, CatBoost, and RF, come from the separation/geometry terms ($\Stk/\sqrt{\Rey}$, $ab/D_e^2$, and the two interaction terms).
\end{enumerate}

\begin{table}[!ht]\centering\caption{LOGO ablation by feature configuration: pooled $R^2$ / pooled RMSE (fixed hyperparameters, constants fitted within folds).}\label{tab:abl}\footnotesize
\begin{tabular}{lccc}\toprule Configuration (no. of inputs) & GBR & XGB & LGBM \\ \midrule
Raw (7) & 0.731 / 0.176 & 0.716 / 0.181 & 0.737 / 0.174 \\
Raw + logs (11) & 0.727 / 0.177 & 0.716 / 0.181 & 0.737 / 0.174 \\
Raw + $\Pi_{VC}$ (8) & 0.729 / 0.176 & 0.715 / 0.181 & 0.735 / 0.174 \\
Raw + separation/geometry terms (12) & 0.730 / 0.176 & 0.735 / 0.174 & 0.750 / 0.169 \\
All features (17) & 0.740 / 0.173 & 0.739 / 0.173 & 0.746 / 0.171 \\
\bottomrule\end{tabular}
\vspace{8pt}

\begin{tabular}{lccc}\toprule Configuration (no. of inputs) & CatBoost & RF & HGB \\ \midrule
Raw (7) & 0.696 / 0.187 & 0.719 / 0.180 & 0.731 / 0.176 \\
Raw + logs (11) & 0.692 / 0.188 & 0.716 / 0.180 & 0.731 / 0.176 \\
Raw + $\Pi_{VC}$ (8) & 0.704 / 0.184 & 0.713 / 0.182 & 0.741 / 0.172 \\
Raw + separation/geometry terms (12) & 0.724 / 0.178 & 0.738 / 0.174 & 0.732 / 0.175 \\
All features (17) & 0.721 / 0.179 & 0.740 / 0.173 & 0.731 / 0.176 \\
\bottomrule\end{tabular}
\end{table}

The paired per-geometry RMSE differences under LOGO have bootstrap intervals that include zero for every model (e.g.\ GBR: $+0.003$, 95\% interval $[-0.013, 0.020]$). The gains seen in Table~\ref{tab:tuned} under tuned grouped splits are therefore not confirmed as statistically significant under the more demanding LOGO protocol. This is an important distinction: the grouped-split gains concern models evaluated on randomly selected subsets of geometries, while LOGO evaluates on every geometry in turn. The correct conclusion is that the engineered features provide modest improvements under repeated geometry-grouped validation, but no statistically significant improvement in leave-one-geometry-out prediction.

All MLP experiments use a \texttt{StandardScaler} fitted exclusively on the training data within each fold, implemented through a scikit-learn pipeline. Under LOGO, the standardized MLP yields pooled $R^2$ values of 0.598 and 0.594 for the raw and 17-feature configurations, respectively, with corresponding pooled RMSE values of 0.215 and 0.216. The macro-RMS RMSE decreases from 0.235 to 0.198, but the paired bootstrap 95\% confidence interval for this difference, $[-0.080, +0.011]$, includes zero. Thus, the point estimate suggests a possible improvement in average per-geometry error, but no statistically significant LOGO improvement is established. The larger MLP gain reported in Table~\ref{tab:fixed} ($R^2$: 0.51 to 0.69) pertains to repeated geometry-grouped splits with fixed hyperparameters and does not transfer to LOGO.

\subsection{Calibration of $\PiVC$}
\label{sec:calibration}

When $(A_\Pi, B_\Pi, C_\Pi)$ are calibrated within training folds (objective function: Eq.~\eqref{eq:objective}), the constants are poorly identified. The damping term collapses ($B_\Pi \approx 10^{-10}$, $C_\Pi$ at its lower bound 0.1 in every fold) and $A_\Pi \approx 10^{5}$, so the values $A_\Pi=0.70$, $B_\Pi=0.30$, $C_\Pi=2.0$ reported in the first submission are not reproducible by out-of-fold fitting. When the Reynolds exponent $m_\Pi$ is also fitted, it converges to $m_\Pi \approx +0.43$ (range 0.40--0.47 across 28 folds), not $-0.5$. This reflects the fact that Stk already contains a factor of Re through its definition, so freeing $m_\Pi$ allows the optimizer to partially compensate for collinearity rather than recovering an independent physical scaling. Freeing $m_\Pi$ does not improve LOGO accuracy (pooled $R^2$ 0.728 versus 0.740 with fixed $m_\Pi=-0.5$). We accordingly describe $\PiVC$ as a data-fitted dimensionless group, not a physically validated descriptor.

\subsection{Generalization to unseen geometries}
\label{sec:logo}

Table~\ref{tab:logo} reports the full LOGO results. For GBR with 17 features, pooled $R^2 = 0.740$ (RMSE 0.173, MAE 0.126), with a weighted mean per-geometry $R^2_g = 0.65 \pm 0.49$ and median 0.75; 23 geometries have $R^2_g > 0.5$ and two have $R^2_g \le 0$. L2GO (60 random pairs) gives pooled $R^2 = 0.737$ (raw: 0.725), with weighted $R^2_g = 0.65 \pm 0.67$ (raw: $0.67 \pm 0.35$). The large standard deviation of $R^2_g$ is mostly a property of the metric itself on small geometries (e.g.\ $R^2_g = -3.9$ for a 9-point geometry with RMSE 0.26), which is why we favour RMSE-based metrics.

Pooled RMSE weights geometries by their number of points, so a geometry with 161 points contributes far more than one with 7. Since the practical goal is prediction for a \emph{new} geometry regardless of how many points it might have, Table~\ref{tab:logo} also reports the macro-RMS of the 28 per-geometry RMSE values ($\sqrt{G^{-1}\sum_g \text{RMSE}_g^2}$). For GBR with 17 features, this macro-RMS is 0.180, compared to 0.173 pooled and 0.165 arithmetic mean. The macro-RMS is consistently slightly higher than the pooled value because the smaller geometries---which receive equal weight in the macro-RMS---tend to have somewhat higher individual RMSE.

\begin{table}[!ht]\centering\caption{LOGO validation (28 folds), all models. Pooled $R^2$, pooled RMSE, macro-RMS RMSE, weighted mean per-geometry $R^2_g$, median $R^2_g$, and number of geometries with $R^2_g \le 0$.}\label{tab:logo}\footnotesize\setlength{\tabcolsep}{3.5pt}\input{t_logo.tex}\end{table}

\begin{figure}[!ht]\centering\includegraphics[width=\textwidth]{fig_logo.png}\caption{LOGO results for GBR (17 features): (a) out-of-geometry parity plot, (b) RMSE per geometry with red markers indicating Mahalanobis-flagged designs, (c) RMSE versus leave-one-out Mahalanobis distance.}\label{fig:logo}\end{figure}

A critical observation concerns the difference between pooled and macro-RMS improvements. GBR's pooled RMSE decreases from 0.176 (raw) to 0.173 (17 features), but the macro-RMS increases from 0.179 to 0.180. The same pattern---pooled improvement but macro-RMS degradation---appears for RF (0.180$\to$0.173 pooled, but 0.180$\to$0.187 macro) and HGB (0.176$\to$0.176 pooled, but 0.175$\to$0.178 macro). In other words, the 17 features improve performance on the larger geometries (which dominate the pooled metric) while slightly degrading average per-geometry performance. This reinforces the conclusion that the benefit of the engineered features is not universal.

\paragraph{Bootstrap of the macro-RMS difference.}
To test this conclusion statistically, we compute the paired bootstrap of the macro-RMS difference
\begin{equation}
\Delta\text{RMSE}_{\text{macro,RMS}} = \text{RMSE}_{\text{macro,RMS}}^{(\text{17 feat.})} - \text{RMSE}_{\text{macro,RMS}}^{(\text{raw})},
\end{equation}
with 2000 replicates resampling geometries (the same 28 geometries in each replicate, paired). The results, reported in Table~\ref{tab:bootstrap}, show that the 95\% bootstrap interval contains zero for every model. This applies to both the macro-RMS difference and the arithmetic-mean difference. The evidence is therefore unambiguous: the engineered features do not produce a statistically significant improvement in leave-one-geometry-out prediction, whether the metric is pooled RMSE, macro-RMS, or arithmetic mean.

\begin{table}[!ht]\centering
\caption{Paired bootstrap (2000 replicates) of the macro-RMS difference $\Delta\text{RMSE}_{\text{macro,RMS}}$ (17 feat.\ $-$ raw) under LOGO. CI: 95\% bootstrap interval. ``CI$\ni$0'': whether the interval contains zero.}
\label{tab:bootstrap}\small
\begin{tabular}{lrrrrl}
\toprule
Model & macro (raw) & macro (17) & $\Delta$ & 95\% CI & CI$\ni$0? \\
\midrule
GBR      & 0.179 & 0.180 & $+0.000$ & $[-0.017, +0.017]$ & yes \\
XGB      & 0.185 & 0.178 & $-0.007$ & $[-0.033, +0.016]$ & yes \\
LGBM     & 0.172 & 0.173 & $+0.002$ & $[-0.013, +0.018]$ & yes \\
CatBoost & 0.188 & 0.179 & $-0.009$ & $[-0.019, +0.001]$ & yes \\
RF       & 0.180 & 0.187 & $+0.007$ & $[-0.009, +0.026]$ & yes \\
HGB      & 0.175 & 0.178 & $+0.003$ & $[-0.015, +0.024]$ & yes \\
MLP      & 0.235 & 0.198 & $-0.036$ & $[-0.080, +0.011]$ & yes \\
\bottomrule
\end{tabular}
\end{table}

\begin{figure}[!ht]\centering\includegraphics[width=\textwidth]{fig_ablation.png}\caption{Ablation results: (a) $R^2$ over 20 grouped splits (error bars: standard deviation), (b) pooled LOGO RMSE for each feature configuration.}\label{fig:abl}\end{figure}

\subsection{Reynolds-number regimes}
Table~\ref{tab:re} gives the LOGO error broken down by Reynolds-number band (edges at $2.3\times10^4$, $5.15\times10^4$, $2.5\times10^5$) and the accuracy when an entire band is withheld. Error is lowest in the highest band (RMSE 0.122; 7 geometries) and 0.17--0.19 elsewhere. The bands differ substantially in composition (4 to 12 geometries, different density-ratio regimes: median $\rho_p/\rho_g$ of 816, 875, 2250, 2250), so differences cannot be attributed to Reynolds number alone. Withholding a whole band is much worse than LOGO ($R^2$ ranging from $-0.14$ for the lowest band to 0.84 for the highest), which is a clear warning that the model should not be used outside the Reynolds range of the training geometries: $2.4\times10^3$--$7.6\times10^5$.

\begin{table}[!ht]\centering\caption{Reynolds-number bands and LOGO error per band.}\label{tab:re}\input{t_re.tex}\end{table}

\subsection{Mahalanobis indicator}
With distances computed against the 27 remaining geometries in each LOGO fold, 6 geometries are flagged. Two of the five geometries with $R^2_g \le 0.5$ are among the flagged set, while other flagged geometries are predicted well (e.g.\ $R^2_g = 0.95$, 0.88). The Spearman correlation between distance and LOGO RMSE is 0.13 ($p = 0.50$); the correlation with $R^2_g$ is $-0.09$ ($p = 0.65$). The AUC for identifying the worst-quartile geometries is 0.62 (bootstrap 95\% interval 0.33--0.89). A $k$-nearest-neighbour distance does no better (Spearman 0.02). The indicator is therefore \emph{not validated} as a predictor of failure; we retain it only as a descriptive measure of how far a design sits from the database centroid in geometry space and explicitly withdraw the practitioner rule proposed in the first submission. (The thresholds quoted earlier, 3.499 and 3.1, were inconsistent; the fold-wise 95th-percentile values reported here are 2.7--2.9.)

\subsection{Prediction intervals}
Residuals are heteroscedastic, as confirmed by a Breusch--Pagan test on LOGO residuals against $\hat\eta$ (LM $= 35.4$, $p = 3\times10^{-9}$). The residual standard deviation peaks near $\hat\eta \approx 0.5$ ($\sigma \approx 0.22$; quadratic coefficients $c_2 = -0.54$, $c_1 = 0.49$, $c_0 = 0.11$). Under LOGO, the empirical coverage of $\hat\eta \pm k\sigma_\eta$ is 67.0\%, 90.1\%, and 97.3\% for $k = 1, 2, 3$ respectively (nominal: 68.3\%, 95.4\%, 99.7\%). The model is therefore adequate near $1\sigma$ and under-covers at the 95\% level. The 95\% quantile of $|\text{res}|/\sigma$ is 2.55, meaning the 95\% interval should use $\pm 2.55\sigma$ ($\approx 1.3\times$ the Gaussian 1.96). This inflation factor is essentially the same for flagged (2.31) and unflagged geometries (2.55), so the earlier recommendation to widen intervals ``by a factor of 2 or higher'' outside a Mahalanobis threshold is unsupported and is withdrawn.

A conformal rescaling with the quantile $q$ computed on the out-of-fold residuals of the other 27 geometries (mean $q = 2.29$) yields pooled coverage of 92.4\% and mean per-geometry coverage of 95.0\% for a 95\% target, with a minimum of 64.5\% for the worst geometry. The 95\% figure represents empirical mean per-geometry coverage within the available database, rather than a formal guarantee for an arbitrary unseen geometry; it depends on the new geometry resembling the population used for calibration.

\subsection{Cut diameter}
$\Stk_{50}$ was extracted for all 73 conditions (median bootstrap coefficient of variation 3.8\%, much smaller than the 12.4\% reported previously, because here each condition has its own grade-efficiency curve rather than one curve per geometry). Table~\ref{tab:stk} and Fig.~\ref{fig:stk} give the Reynolds exponent. The pooled value ($-0.46$) and the geometry-adjusted values ($-0.55$ and $-0.52$) are all compatible with $-0.5$, but the intervals are wide (up to $\pm 0.2$). The fixed-effects estimate rests on only 17 geometries that have more than one Reynolds number (median max/min Re ratio 1.7), and the unweighted fixed-effects estimate is $-0.66$. Because $\ln\Rey$ is 98.7\% explained by the geometry mean, the pooled slope mostly reflects between-geometry differences rather than a within-geometry physical scaling. We therefore report the exponent as consistent with---but not a validation of---the Barth-type scaling $\Stk_{50} \propto \Rey^{-1/2}$, and we make no physical interpretation of the 7.9\% departure. All uncertainties are conditional on the assumed noise $\sigma_\epsilon = 0.025$ and describe the sensitivity of the extraction procedure, not measured experimental error. Table~\ref{tab:sens} shows that the median coefficient of variation scales from 1.5\% to 8.2\% for $\sigma_\epsilon = 0.01$ to 0.05, while the estimates of $\alpha$ remain between $-0.47$ and $-0.57$ (pooled and covariate-adjusted) and $-0.50$ to $-0.57$ (fixed effects).

\begin{table}[!ht]\centering\caption{Weighted regression of $\ln\Stk_{50}$ on $\ln\Rey$.}\label{tab:stk}\input{t_stk.tex}\end{table}
\begin{table}[!ht]\centering\caption{Sensitivity of the Stk$_{50}$ analysis to the assumed noise level.}\label{tab:sens}\footnotesize\setlength{\tabcolsep}{5pt}\input{t_sens.tex}\end{table}
\begin{figure}[!ht]\centering\includegraphics[width=0.55\textwidth]{fig_stk50.png}\caption{$\Stk_{50}$ versus Re for 73 conditions (colour: geometry; bars: $\pm 2$ bootstrap s.d.), pooled fit and slope $-0.5$ reference line.}\label{fig:stk}\end{figure}

% =====================================================================
%  DISCUSSION AND LIMITATIONS
% =====================================================================
\section{Discussion and limitations}
\label{sec:related}

The results of this study carry implications that extend beyond the specific database analysed. Adaptive optimization and control frameworks for hydrocyclones \cite{rineng}, flow-field surrogates trained on CFD data \cite{mineng,e2024}, and hybrid machine-learning/mechanistic models \cite{seppur, e2025} are typically validated on data from a single design family. Our findings suggest that such models should be reported with geometry-grouped validation in addition to random splitting, because random validation overstates accuracy by roughly 0.2 in $R^2$ in this database---and the overstatement could be larger in databases with fewer geometries or more within-geometry points.

\label{sec:limits}
Several limitations should be noted explicitly:
\begin{enumerate}[label=(\roman*)]
\item Only 28 geometries are available, each defined by three ratios. Body and cone heights and inlet type are not recorded, so physically different cyclones can share a single geometry identity, and the model may partly learn differences between experimental sources or undocumented configurations rather than geometry itself. Inter-source bias is not identifiable in this database.
\item No high-solids-loading data, non-standard inlet configurations, or elevated-temperature measurements are included.
\item The noise level $\sigma_\epsilon$ used in the bootstrap is an assumption, not a measured quantity.
\item The L2GO analysis uses a sample of 60 of 378 possible pairs, and the tuned protocol uses 10 splits, for computational tractability.
\item The pressure-drop dataset is not available in the supplementary material of \cite{bregolin}, so pressure-drop results and the closed-form fit of the first submission are removed rather than reported without a reproducible protocol.
\item The Mahalanobis analysis operates on only 28 points and therefore has low statistical power. Absence of evidence of an association is not evidence of absence, but the indicator cannot be recommended in its current form.
\end{enumerate}

Our results quantify statistical transfer between the 28 geometry groups present in the database, not extrapolation to a new industrial design that differs in unrecorded features. This is a fundamental limitation of any database with incomplete geometric descriptions.

% =====================================================================
%  CONCLUSIONS
% =====================================================================
\section{Conclusions}

This study set out to provide an honest, reproducible assessment of what machine-learning models can and cannot achieve when applied to cyclone collection-efficiency prediction across geometries. The findings can be summarized as follows.

Random point-wise validation, the protocol used in virtually all prior studies on this database, yields $R^2 \approx 0.95$---a number that looks impressive but measures interpolation along known grade-efficiency curves rather than prediction of unseen designs. When whole geometries are withheld from training, $R^2$ drops to 0.73--0.75 (RMSE $\approx 0.17$), with a spread of 0.10--0.14 across repeated splits. This gap of roughly 0.2 in $R^2$ is the central quantitative result of the paper, and it implies that published accuracy figures should be interpreted with caution when the intended use involves a new geometry.

The six tree-based learners we tested exhibit similar average performance under geometry-grouped validation, with split-to-split variability substantially larger than the differences between their mean scores. No clear ranking can therefore be established on this database. All of them comfortably outperform the Lapple ($R^2 \approx 0.29$) and Stokes-only ($R^2 \approx 0.44$) baselines.

The 17 engineered features---including a vortical-capture group $\PiVC$ calibrated within training folds---yield small, consistent gains under tuned hyperparameters on grouped splits (RMSE reduction of 0.001--0.011) and a substantial gain for the neural network under grouped splits with fixed hyperparameters ($R^2$: 0.51$\to$0.69). However, neither gain transfers to LOGO validation: the paired bootstrap 95\% interval of the macro-RMS difference contains zero for every model. The constants of $\PiVC$ are poorly identified out-of-fold, and a freely fitted Reynolds exponent converges to $+0.43$ rather than the assumed $-0.5$, reflecting the intrinsic collinearity between Stk and Re. We therefore describe $\PiVC$ as a data-fitted dimensionless group rather than a validated physical descriptor, and we note that the separation/geometry terms ($\Stk/\sqrt{\Rey}$, $ab/D_e^2$) contribute more to the modest gains observed than does $\PiVC$ itself.

The Mahalanobis extrapolation indicator is not validated as a predictor of LOGO failure (Spearman $\rho = 0.13$, $p = 0.50$), and we withdraw the practitioner rule proposed earlier. The heteroscedastic interval model under-covers at the 95\% level and requires an inflation factor of approximately 1.3; a geometry-calibrated conformal rescaling achieves 95\% mean per-geometry coverage but only 64.5\% for the worst geometry, so this represents empirical mean coverage within the available database, not a formal guarantee for an arbitrary unseen geometry. The cut-diameter exponent ($-0.46$ to $-0.55$) is compatible with the Barth-type $\Rey^{-1/2}$ scaling but cannot be disentangled from geometry effects in this database, where Re and geometry are 98.7\% collinear.

Taken together, these results do not constitute a failure of machine learning for cyclone efficiency prediction---the models clearly outperform semi-empirical correlations---but they do establish a conservative, honest baseline against which future improvements can be measured, and they identify the specific conditions under which current data-driven models should and should not be trusted. The database, the code, and all results are publicly available to support further work.

\section*{Data and code availability}
The efficiency data are from the supplementary material of Bregolin et al.\ \cite{bregolin}. All code (feature construction, fold-wise calibration, nested validation protocols, bootstrap analysis, table and figure generation), the full set of numerical results (CSV and JSON), and all figures are available at\\
\texttt{https://github.com/Ahmedalgeria778/cyclone-efficiency-ml-validation}.

\section*{CRediT}
A.~Hanifi: conceptualization, methodology, software, investigation, writing (original draft). A.~Beniaiche: supervision, writing (review \& editing), revision analyses. A.~Belkalouche, S.~Ghodbane, C.~Belkacemi, M.~Belkadi: validation, resources, review \& editing.

\appendix
\section{Per-geometry characteristics}
\input{t_geom.tex}

\section{Bootstrap macro-RMS details}
The paired bootstrap procedure resamples the 28 geometries (with replacement, same set in each replicate for both raw and 17-feature conditions) and recomputes the macro-RMS difference. Table~\ref{tab:bootstrap_app} gives the full results, including both the macro-RMS and the arithmetic-mean differences.

\begin{table}[!ht]\centering
\caption{Full paired bootstrap results (2000 replicates): macro-RMS and arithmetic-mean differences.}
\label{tab:bootstrap_app}\footnotesize\setlength{\tabcolsep}{3pt}
\begin{tabular}{lrrrrrrrr}
\toprule
& \multicolumn{4}{c}{Macro-RMS} & \multicolumn{4}{c}{Arithmetic mean} \\
\cmidrule(lr){2-5}\cmidrule(lr){6-9}
Model & raw & 17 feat. & $\Delta$ & 95\% CI & raw & 17 feat. & $\Delta$ & 95\% CI \\
\midrule
GBR      & 0.179 & 0.180 & $+0.000$ & $[-0.017, +0.017]$ & 0.162 & 0.165 & $+0.003$ & $[-0.014, +0.019]$ \\
XGB      & 0.185 & 0.178 & $-0.007$ & $[-0.033, +0.016]$ & 0.164 & 0.163 & $-0.001$ & $[-0.020, +0.016]$ \\
LGBM     & 0.172 & 0.173 & $+0.002$ & $[-0.013, +0.018]$ & 0.158 & 0.159 & $+0.001$ & $[-0.013, +0.016]$ \\
CatBoost & 0.188 & 0.179 & $-0.009$ & $[-0.019, +0.001]$ & 0.174 & 0.168 & $-0.006$ & $[-0.017, +0.004]$ \\
RF       & 0.180 & 0.187 & $+0.007$ & $[-0.009, +0.026]$ & 0.164 & 0.169 & $+0.005$ & $[-0.009, +0.020]$ \\
HGB      & 0.175 & 0.178 & $+0.003$ & $[-0.015, +0.024]$ & 0.160 & 0.163 & $+0.003$ & $[-0.013, +0.021]$ \\
MLP      & 0.235 & 0.198 & $-0.036$ & $[-0.080, +0.011]$ & 0.207 & 0.184 & $-0.023$ & $[-0.060, +0.016]$ \\
\bottomrule
\end{tabular}
\end{table}

\clearpage
\begin{thebibliography}{99}\small
\bibitem{hoffmann}A.C. Hoffmann, L.E. Stein, Gas Cyclones and Swirl Tubes, 2nd ed., Springer, Berlin, 2007.
\bibitem{bregolin}E. Bregolin, P. Danieli, M. Masi, Collection efficiency of cyclone separators: Comparison between new machine learning-based models and semi-empirical approaches, Waste 2 (3) (2024) 240--257. doi:10.3390/waste2030014
\bibitem{rineng}Dianyu E, Yuhao Zhang, Cong Tan, Hongwei Hu, Jiaxin Cui, Chengfang Yuan, Zongyan Zhou, Caibin Wu, Ruiping Zou, Shibo Kuang, Aibing Yu, An adaptive multi-objective optimization and control framework for hydrocyclone operation under dynamic particle-size distributions, Results in Engineering 29 (2026) 109026. doi:10.1016/j.rineng.2026.109026
\bibitem{mineng}Dianyu E, Yuhao Zhang, Hongwei Hu, Guangtai Xu, Lin Liu, Jiaxin Cui, Cong Tan, Ruiping Zou, Shibo Kuang, CFD data-driven CNN-LSTM for fast prediction of instantaneous flow characteristics in hydrocyclones, Minerals Engineering 235 (2026) 109814. doi:10.1016/j.mineng.2025.109814
\bibitem{seppur}Wenbo Wu, Yuhao Zhang, Cong Tan, Jiaxin Cui, Chengfang Yuan, Zongyan Zhou, Caibin Wu, Ruiping Zou, Dianyu E, Shibo Kuang, A combined machine learning and mechanistic model for quickly predicting three-dimensional inner states of hydrocyclones, Separation and Purification Technology 379 (2025) 135046. doi:10.1016/j.seppur.2025.135046
\bibitem{e2024}Dianyu E, Guangtai Xu, Jiaxin Cui, Qing Ye, Cong Tan, Ruiping Zou, Aibing Yu, Shibo Kuang, Prediction of instantaneous flow characteristics of hydrocyclone with long short-term memory network based on computational fluid dynamics data, Powder Technology 439 (2024) 119668. doi:10.1016/j.powtec.2024.119668
\bibitem{e2025}Dianyu E, Hongwei Hu, Qing Ye, Cong Tan, Jiaxin Cui, Zongyan Zhou, Ruiping Zou, Aibing Yu, Shibo Kuang, Multi-objective optimization of hydrocyclones inlet configurations for improving separation performance, Powder Technology 456 (2025) 120772. doi:10.1016/j.powtec.2025.120772
\bibitem{lapple}C.E. Lapple, Processes use many collector types, Chemical Engineering 58 (5) (1951) 144--151.
\bibitem{barth}W. Barth, Berechnung und Auslegung von Zyklonabscheidern auf Grund neuerer Untersuchungen, Brennstoff-W\"arme-Kraft 8 (1956) 1--9.
\bibitem{muschel}E. Muschelknautz, M. Trefz, Design and calculation of cyclones, in: VDI Heat Atlas, VDI-Verlag, D\"usseldorf, 1990.
\bibitem{hoffmann92}A.C. Hoffmann, A. van Santen, R.W.K. Allen, R. Clift, Effects of geometry and solid loading on the performance of gas cyclones, Powder Technology 70 (1992) 83--91.
\bibitem{dirgo}J. Dirgo, D. Leith, Cyclone collection efficiency: comparison of experimental results with theoretical predictions, Aerosol Science and Technology 4 (1985) 401--415.
\bibitem{elsayed}K. Elsayed, C. Lacor, The effect of cyclone inlet dimensions on the flow pattern and performance, Applied Mathematical Modelling 35 (4) (2011) 1952--1968. doi:10.1016/j.apm.2010.11.007
\bibitem{zhao}B. Zhao, Y. Su, Artificial neural network-based modeling of pressure drop coefficient for cyclone separators, Chemical Engineering Research and Design 88 (5--6) (2010) 606--613. doi:10.1016/j.cherd.2009.11.010
\bibitem{karniadakis}G.E. Karniadakis, I.G. Kevrekidis, L. Lu, P. Perdikaris, S. Wang, L. Yang, Physics-informed machine learning, Nature Reviews Physics 3 (2021) 422--440.
\end{thebibliography}
\end{document}
